% !TEX program = xelatex
\documentclass[aspectratio=169,10pt]{beamer}

\usepackage[T1]{fontenc}
\usepackage[default]{sourcesanspro}
\usepackage{graphicx}
\usepackage{microtype}

% Monochrome palette for presentation elements. Source figures retain their
% original appearance so that the published evidence remains visible.
\definecolor{Canvas}{HTML}{F5F5F5}
\definecolor{Paper}{HTML}{FFFFFF}
\definecolor{Ink}{HTML}{1E1E1E}
\definecolor{DarkGray}{HTML}{383838}
\definecolor{Muted}{HTML}{666666}
\definecolor{Line}{HTML}{CCCCCC}

\setbeamercolor{normal text}{fg=Ink,bg=Canvas}
\setbeamercolor{frametitle}{fg=Ink,bg=Canvas}
\setbeamercolor{structure}{fg=DarkGray}
\setbeamercolor{itemize item}{fg=DarkGray}
\setbeamercolor{itemize subitem}{fg=Muted}
\setbeamerfont{frametitle}{size=\fontsize{15}{17}\selectfont,series=\bfseries}
\setbeamerfont{normal text}{size=\fontsize{9.5}{11.7}\selectfont}
\setbeamertemplate{navigation symbols}{}
\setbeamertemplate{footline}{}
\setbeamertemplate{headline}{}
\setbeamersize{text margin left=8mm,text margin right=8mm}
\setbeamertemplate{frametitle}[default][left]
\setbeamertemplate{itemize item}{\color{DarkGray}\raisebox{.18ex}{\scriptsize\(\bullet\)}}
\setbeamertemplate{itemize subitem}{\color{Muted}\raisebox{.18ex}{\scriptsize\(\circ\)}}
\setlength{\leftmargini}{1.55em}
\setlength{\leftmarginii}{1.35em}

\newcommand{\source}[1]{{\color{Muted}\fontsize{6.4}{7.2}\selectfont #1\par}}
\newcommand{\term}[1]{{\color{Ink}#1}}
\newcommand{\metric}[1]{{\color{DarkGray}#1}}
\newcommand{\paperimage}[2]{%
  \begingroup%
  \setlength{\fboxsep}{1.2pt}%
  \setlength{\fboxrule}{.35pt}%
  \fcolorbox{Line}{Paper}{\includegraphics[#1]{#2}}%
  \endgroup%
}
\newenvironment{tightitems}{%
  \begin{itemize}\setlength{\itemsep}{.52em}\setlength{\topsep}{.15em}\setlength{\parskip}{0pt}%
}{\end{itemize}}
\newenvironment{compactitems}{%
  \begin{itemize}\setlength{\itemsep}{.25em}\setlength{\topsep}{.1em}\setlength{\parskip}{0pt}%
}{\end{itemize}}

\hypersetup{pdftitle={V\&V setups, activities and objectives for nanosatellite EPS},pdfauthor={}}

\begin{document}

% -----------------------------------------------------------------------------
\begin{frame}[plain]
  \vspace*{17mm}
  {\raggedright\fontsize{28}{31}\selectfont\bfseries
    V\&V setups, activities, and objectives\\[-1mm]
    for nanosatellite EPS\par}
  \vfill
\end{frame}

% -----------------------------------------------------------------------------
\begin{frame}{Categories of setup context and intended evidence}
  \begin{columns}[T,onlytextwidth]
    \begin{column}{.485\textwidth}
      \begin{tightitems}
        \item \term{Executable and analytical models}. Software represents orbit, loads, generation, battery state, requirements, and failures. It supports sizing, compliance checks, and reliability analysis.
        \item \term{Component bench}. One physical item is measured under controlled conditions. The resulting curves supply model parameters.
        \item \term{Real-time digital twin / HIL}. Simulated and real components exchange signals at test speed. This exposes transients and controller behavior.
      \end{tightitems}
    \end{column}
    \begin{column}{.485\textwidth}
      \begin{tightitems}
        \item \term{FlatSat / EGSE}. Subsystems are wired with ground equipment before full spacecraft integration. This checks interfaces, commands, and data paths.
        \item \term{Environmental}. A flight or qualification article faces thermal or vacuum conditions. This checks survival, battery safety, and retained function.
        \item \term{Flight / telemetry}. The operational spacecraft executes schedules and returns measurements. This assesses depletion risk and mission outcome.
      \end{tightitems}
    \end{column}
  \end{columns}
\end{frame}

% -----------------------------------------------------------------------------
\begin{frame}{V\&V questions supported by the setups}
  \begin{columns}[T,onlytextwidth]
    \begin{column}{.485\textwidth}
      \begin{tightitems}
        \item \term{Desktop and analytical models}. Worst-case runs, design of experiments, benchmarks, and failure analyses support sizing, compliance, and risk reduction.
        \item \term{Component bench}. I--V, P--V, and discharge curves support empirical solar and battery models.
        \item \term{Real-time twin, CHIL, and PHIL}. Physical transients and control changes support model credibility and controller robustness.
        \item \term{FlatSat, EGSE, and solar simulator}. Commands, data, resets, load shedding, and illumination checks support interface correctness and generation capacity.
      \end{tightitems}
    \end{column}
    \begin{column}{.485\textwidth}
      \begin{tightitems}
        \item \term{Thermal vacuum}. Hot and cold cycles with health checks support battery safety and readiness.
        \item \term{Qualification telemetry}. Injected faults and residual checks support early fault isolation.
        \item \term{Flight telemetry loop}. Scheduled loads and measured battery response support depletion avoidance and mission yield.
      \end{tightitems}
    \end{column}
  \end{columns}
\end{frame}

% -----------------------------------------------------------------------------
\begin{frame}{SysML traceability for V\&V coverage and verification status}
  \begin{columns}[T,onlytextwidth]
    \begin{column}{.48\textwidth}
      \centering
      \paperimage{width=.92\linewidth}{assets/citation-1691-test-relation.jpg}\par\vspace{2mm}
      \paperimage{width=.92\linewidth}{assets/citation-1691-status-table.jpg}
    \end{column}
    \begin{column}{.48\textwidth}
      \begin{tightitems}
        \item The template links EPS test activities to requirements and records verification status.
        \item Its visible requirements cover power, energy storage, fault protection, load, charging, temperature, and state of charge.
        \item The source marks the mass example as verified and makes the remaining test links visible.
        \item The template supports traceability. Verified test evidence completes EPS coverage.
      \end{tightitems}
    \end{column}
  \end{columns}
  \vfill
  \source{Source figures: \emph{A reference architecture for rapid cubesat development}, Figures 45--46, pp. 88--89.}
\end{frame}

% -----------------------------------------------------------------------------
\begin{frame}{Mission simulations expose margin and transient depletion}
  \begin{columns}[T,onlytextwidth]
    \begin{column}{.48\textwidth}
      \centering
      \paperimage{width=.97\linewidth}{assets/review-0289-stable-budget.jpg}\par\vspace{1mm}
      {\color{Muted}\fontsize{6.8}{7.5}\selectfont 48-HOUR WORST-ILLUMINATION CASE\par}
    \end{column}
    \begin{column}{.48\textwidth}
      \centering
      \paperimage{width=.97\linewidth}{assets/review-1231-depletion.jpg}\par\vspace{1mm}
      {\color{Muted}\fontsize{6.8}{7.5}\selectfont SEVEN-DAY NETSAT BATTERY STATE\par}
    \end{column}
  \end{columns}
  \vspace{1mm}
  \begin{compactitems}
    \item The CNES IDM-CIC, IDM-Power, Simu-CIC, and Opalis chain varied array, battery, regulator, and architecture parameters. The selected design retained a \metric{85\%} minimum state of charge.
    \item ESTNeT coupled orbit-control behavior, solar orientation, loads, and battery state. It found repeated depletion on one satellite that the positive average power budget missed.
  \end{compactitems}
  \vfill
  \source{Source figures: \emph{Concurrent Engineering Process for Electrical Architecture of Nanosatellites Preliminary Phases}, Figure 10, pp. 2--6; \emph{CubeSat energy modelling for improved mission planning and operations}, Figure 3, pp. 5--7.}
\end{frame}

% -----------------------------------------------------------------------------
\begin{frame}{Requirement-linked simulations test explicit energy margins}
  \begin{columns}[T,onlytextwidth]
    \begin{column}{.48\textwidth}
      \centering
      \paperimage{height=3.20cm}{assets/review-0499-soc.jpg}\par\vspace{1mm}
      {\color{Muted}\fontsize{6.8}{7.5}\selectfont 3DEXPERIENCE / DYMOLA STATE OF CHARGE\par}
    \end{column}
    \begin{column}{.48\textwidth}
      \centering
      \paperimage{height=3.20cm}{assets/review-1563-solar-sizing.jpg}\par\vspace{1mm}
      {\color{Muted}\fontsize{6.8}{7.5}\selectfont RAX SOLAR-AREA REQUIREMENT TEST\par}
    \end{column}
  \end{columns}
  \vspace{1mm}
  \begin{compactitems}
    \item A requirements-in-loop workflow ran hundreds of battery configurations against a \metric{60\%} state-of-charge floor and recharge objective. The reported case remained above \metric{90\%} for 300 minutes.
    \item A SysML, STK, MATLAB, and ModelCenter workflow checked a \metric{92,000 J} minimum-energy requirement. Nominal and half-size arrays passed. The quarter-size array crossed the limit.
  \end{compactitems}
  \vfill
  \source{Source figures: \emph{An Implementation of Integrated Approach in Product Life-Cycle Management Tool to Ensure Requirements-In-Loop During Complex Product Development}, Figure 6, pp. 5--9; \emph{Integrated model-based systems engineering applied to the Simulation of a CubeSat mission}, Figure 16, pp. 6--12.}
\end{frame}

% -----------------------------------------------------------------------------
\begin{frame}{Reusable V\&V platform: models, test control, and physical equipment}
  \begin{columns}[T,onlytextwidth]
    \begin{column}{.58\textwidth}
      \centering
      \paperimage{height=4.85cm}{assets/review-1265-platform.jpg}
    \end{column}
    \begin{column}{.38\textwidth}
      \begin{compactitems}
        \item Requirements, models, and budgets feed test control and instrumentation.
        \item Available assets include automated test equipment, HIL, the device under test, and thermal cycling.
        \item The electrical bench can include power supplies, a DMM, a power meter, an oscilloscope, and standard instrument links.
        \item The published hardware demonstration focused on communications. The platform supports future EPS testing.
      \end{compactitems}
    \end{column}
  \end{columns}
  \vfill
  \source{Source figure: \emph{Rapid, automated, test, verification and validation for the CubeSats}, Figure 2, pp. 7--15.}
\end{frame}

% -----------------------------------------------------------------------------
\begin{frame}{Test architecture scales from prototype to flight article}
  \begin{columns}[T,onlytextwidth]
    \begin{column}{.53\textwidth}
      \centering
      \paperimage{width=.96\linewidth}{assets/review-0239-test-system.jpg}
    \end{column}
    \begin{column}{.43\textwidth}
      {\fontsize{8.2}{9.5}\selectfont
      \begin{compactitems}
        \item The model connects the plan, equipment, instrumentation, procedure, execution, and results.
        \item EPS objectives include battery performance, solar capacity, heat, battery heaters, vibration, EMI, and telemetry power.
        \item Proxy procedures gave relative durations of \metric{1.00}, \metric{1.69}, and \metric{2.78} for prototype, qualification, and flight models.
        \item The solar procedure used an AMSAT CubeSat Simulator and a solar energy simulator. The broader matrix remained proposed.
      \end{compactitems}}
    \end{column}
  \end{columns}
  \vfill
  \source{Source: \emph{A Scalable System Architecture for CubeSat Test \& Evaluation for Enhanced Mission Success}, Figures 25 and 47 and Appendix A, pp. 86--87, 112--113, and 142--155.}
\end{frame}

% -----------------------------------------------------------------------------
\begin{frame}{Solar and battery characterization for the digital twin}
  \begin{columns}[T,onlytextwidth]
    \begin{column}{.48\textwidth}
      \centering
      \paperimage{height=3.48cm}{assets/review-0573-sas-curves.jpg}\par\vspace{1mm}
      {\color{Muted}\fontsize{6.8}{7.5}\selectfont SOLAR-ARRAY CHARACTERIZATION\par}
    \end{column}
    \begin{column}{.48\textwidth}
      \centering
      \paperimage{height=3.48cm}{assets/review-0573-battery-curve.jpg}\par\vspace{1mm}
      {\color{Muted}\fontsize{6.8}{7.5}\selectfont 2S2P BATTERY CHARACTERIZATION\par}
    \end{column}
  \end{columns}
  \vspace{1mm}
  \begin{compactitems}
    \item The solar simulator measured an eight-cell solar-array section at AM0 irradiance and 85 \(^\circ\)C. Its I--V and P--V curves parameterized the model.
    \item A 2S2P Samsung 18650 battery was discharged at 520 mA and 25 \(^\circ\)C. A current step identified series resistance.
  \end{compactitems}
  \vfill
  \source{Source figures: \emph{Implementation of a 6U CubeSat Electrical Power System Digital Twin}, Figures 2 and 7, pp. 3 and 7.}
\end{frame}

% -----------------------------------------------------------------------------
\begin{frame}{Hardware coupling to a 500 ns real-time EPS twin}
  \begin{columns}[T,onlytextwidth]
    \begin{column}{.30\textwidth}
      \centering
      \paperimage{height=3.30cm}{assets/review-0573-setup.jpg}\par\vspace{1mm}
      {\color{Muted}\fontsize{6.6}{7.2}\selectfont HIL404 AND INSTRUMENTS\par}
    \end{column}
    \begin{column}{.35\textwidth}
      \centering
      \paperimage{height=3.30cm}{assets/review-0573-physical-eps.jpg}\par\vspace{1mm}
      {\color{Muted}\fontsize{6.6}{7.2}\selectfont PHYSICAL EPS\par}
    \end{column}
    \begin{column}{.27\textwidth}
      \centering
      \paperimage{height=3.30cm}{assets/review-0573-cubesat-sas.jpg}\par\vspace{1mm}
      {\color{Muted}\fontsize{6.6}{7.2}\selectfont 6U WITH TWO SAS\par}
    \end{column}
  \end{columns}
  \vspace{1mm}
  \begin{compactitems}
    \item A Typhoon HIL404 ran a twin of one solar-array section, regulator, battery, PDU, and loads with a \metric{500 ns} simulation step.
    \item The reported \metric{1.65 \(\mu\)s} loopback latency connected the physical and digital parts during MPPT, end-of-charge regulation, and a load transient.
  \end{compactitems}
  \vfill
  \source{Source figures: \emph{Implementation of a 6U CubeSat Electrical Power System Digital Twin}, Figure 8, pp. 8--10.}
\end{frame}

% -----------------------------------------------------------------------------
\begin{frame}{Transient EPS response to an 18 W load step}
  \begin{columns}[T,onlytextwidth]
    \begin{column}{.58\textwidth}
      \centering
      \paperimage{height=4.92cm}{assets/review-0573-dt-validation.jpg}
    \end{column}
    \begin{column}{.38\textwidth}
      \begin{tightitems}
        \item At end of charge, solar generation exceeded load and recharged the battery.
        \item Adding an \metric{18 W} load at 2 s created a negative balance. The battery discharged and MPPT engaged.
        \item Removing the load at 8 s returned the EPS to end-of-charge regulation.
        \item The source reports almost identical physical and digital waveforms across this sequence.
      \end{tightitems}
    \end{column}
  \end{columns}
  \vfill
  \source{Source figure: \emph{Implementation of a 6U CubeSat Electrical Power System Digital Twin}, Figure 9, pp. 9--10.}
\end{frame}

% -----------------------------------------------------------------------------
\begin{frame}{CHIL MPPT tracking with a physical controller}
  \begin{columns}[T,onlytextwidth]
    \begin{column}{.48\textwidth}
      \centering
      \paperimage{height=3.55cm}{assets/review-0573-chil-diagram.jpg}\par\vspace{1mm}
      {\color{Muted}\fontsize{6.8}{7.5}\selectfont PUBLISHED CHIL SETUP\par}
    \end{column}
    \begin{column}{.48\textwidth}
      \centering
      \paperimage{height=3.55cm}{assets/review-0573-chil-result.jpg}\par\vspace{1mm}
      {\color{Muted}\fontsize{6.8}{7.5}\selectfont PUBLISHED CHIL RESPONSE\par}
    \end{column}
  \end{columns}
  \vspace{1mm}
  \begin{compactitems}
    \item The physical control circuit replaces the simulated regulator controller. The solar array, converter, battery, and loads remain in the real-time simulation.
    \item A solar-cell temperature sweep from \metric{--50 \(^\circ\)C to +100 \(^\circ\)C} exercises MPPT tracking over expected conditions.
  \end{compactitems}
  \vfill
  \source{Source figures: \emph{Implementation of a 6U CubeSat Electrical Power System Digital Twin}, Figures 10--11, pp. 9--11.}
\end{frame}

% -----------------------------------------------------------------------------
\begin{frame}{PHIL evaluation of a digital MPPT controller and physical power stage}
  \begin{columns}[T,onlytextwidth]
    \begin{column}{.48\textwidth}
      \centering
      \paperimage{height=3.55cm}{assets/review-0573-phil-diagram.jpg}\par\vspace{1mm}
      {\color{Muted}\fontsize{6.8}{7.5}\selectfont PUBLISHED PHIL SETUP\par}
    \end{column}
    \begin{column}{.48\textwidth}
      \centering
      \paperimage{height=3.55cm}{assets/review-0573-phil-result.jpg}\par\vspace{1mm}
      {\color{Muted}\fontsize{6.8}{7.5}\selectfont PUBLISHED PHIL RESPONSE\par}
    \end{column}
  \end{columns}
  \vspace{1mm}
  \begin{compactitems}
    \item The real-time simulator runs a digital perturb-and-observe MPPT controller while the converter, battery, and loads remain physical.
    \item The observed operating point oscillated around the approximately \metric{15 V} maximum-power point, supporting assessment of a control alternative.
  \end{compactitems}
  \vfill
  \source{Source figures: \emph{Implementation of a 6U CubeSat Electrical Power System Digital Twin}, Figures 12--13, pp. 10--12.}
\end{frame}

% -----------------------------------------------------------------------------
\begin{frame}{Thermal-vacuum functional testing of the EPS flight article}
  \begin{columns}[T,onlytextwidth]
    \begin{column}{.31\textwidth}
      \centering
      \paperimage{height=2.30cm}{assets/review-0532-tvac-setup.jpg}\par\vspace{1.5mm}
      \paperimage{height=1.70cm}{assets/review-0532-tc-placement.jpg}
    \end{column}
    \begin{column}{.65\textwidth}
      \centering
      \paperimage{width=.95\linewidth}{assets/review-0532-tvac-sequence.jpg}
    \end{column}
  \end{columns}
  \vspace{1mm}
  \begin{compactitems}
    \item The 3Cat-4 flight article entered a thermal-vacuum chamber with twelve thermocouples, including sensors on the batteries and EPS board.
    \item The campaign covered hot and cold conditions, operational and non-operational phases, charging, thermal balance, and reduced functional tests.
  \end{compactitems}
  \vfill
  \source{Source figures: \emph{Contributions to eco-friendly satellite lean design for a sustainable space environment}, Figures 8.8, 8.10, and 8.12, pp. 203--217.}
\end{frame}

% -----------------------------------------------------------------------------
\begin{frame}{Cold battery finding, EPS configuration update, and model refinement}
  \begin{columns}[T,onlytextwidth]
    \begin{column}{.58\textwidth}
      \centering
      \paperimage{height=4.98cm}{assets/review-0532-cold-balance.jpg}
    \end{column}
    \begin{column}{.38\textwidth}
      \begin{tightitems}
        \item A cold switch-on produced a \term{Cold Batteries} error. Heater activation then raised EPS and battery temperatures quickly.
        \item The finding triggered an EPS configuration update.
        \item Thermal Desktop predictions and test measurements refined conductance, material, and radiative properties.
        \item Pearson correlations were \metric{0.7750} for hot balance and \metric{0.78904} for cold balance.
      \end{tightitems}
    \end{column}
  \end{columns}
  \vfill
  \source{Source figure: \emph{Contributions to eco-friendly satellite lean design for a sustainable space environment}, Figure 8.15, pp. 212--219. The overall campaign met its requirements.}
\end{frame}

% -----------------------------------------------------------------------------
\begin{frame}{Hot and cold balance data correlate the thermal model}
  \begin{columns}[T,onlytextwidth]
    \begin{column}{.48\textwidth}
      \centering
      \paperimage{height=3.50cm}{assets/review-0532-hot-correlation.jpg}\par\vspace{1mm}
      {\color{Muted}\fontsize{6.8}{7.5}\selectfont HOT BALANCE: TEST AND MODEL\par}
    \end{column}
    \begin{column}{.48\textwidth}
      \centering
      \paperimage{height=3.50cm}{assets/review-0532-cold-correlation.jpg}\par\vspace{1mm}
      {\color{Muted}\fontsize{6.8}{7.5}\selectfont COLD BALANCE: TEST AND MODEL\par}
    \end{column}
  \end{columns}
  \vspace{1mm}
  \begin{compactitems}
    \item Hot-balance correlation was \metric{0.7750}, with a maximum reported difference of \metric{5.10~\(^\circ\)C} at the COMMS--ADCS thermocouple.
    \item Cold-balance correlation was \metric{0.78904}, with a maximum reported difference of \metric{6.34~\(^\circ\)C} at the COMMS MCU sensor.
    \item The comparison supports refinement of interface conductance, material properties, and radiative properties.
  \end{compactitems}
  \vfill
  \source{Source figures: \emph{Contributions to eco-friendly satellite lean design for a sustainable space environment}, Figures 8.17--8.18, pp. 217--219.}
\end{frame}

% -----------------------------------------------------------------------------
\begin{frame}{Modeled EPS functional and thermal-vacuum verification paths}
  \begin{columns}[T,onlytextwidth]
    \begin{column}{.48\textwidth}
      \centering
      \paperimage{height=3.45cm}{assets/citation-1572-functional-path.jpg}\par\vspace{1mm}
      {\color{Muted}\fontsize{6.8}{7.5}\selectfont BATTERY-BYPASS FUNCTIONAL PATH\par}
    \end{column}
    \begin{column}{.48\textwidth}
      \centering
      \paperimage{height=3.45cm}{assets/citation-1572-tvac-model.jpg}\par\vspace{1mm}
      {\color{Muted}\fontsize{6.8}{7.5}\selectfont TVAC AND ELECTRICAL-TEST INTERFACES\par}
    \end{column}
  \end{columns}
  \vspace{1mm}
  \begin{compactitems}
    \item The AESP-14 functional path joins EPS, TT\&C, a battery-bypass supply, JTAG, EGSE, a PC, a portable ground station, and mechanical support.
    \item The TVAC model adds the chamber, thermocouples, measurement and acquisition systems, EGSE, and power cabling that bypasses the batteries.
    \item Both views define functions, physical elements, interfaces, requirements, and discrepancy impact for test planning.
  \end{compactitems}
  \vfill
  \source{Source figures: \emph{A conceptual MBSE framework for satellite AIT planning}, Figures 6.31 and 6.35, pp. 153--167.}
\end{frame}

% -----------------------------------------------------------------------------
\begin{frame}{DESCENT test architecture, FlatSat, and compressed sequence}
  \begin{columns}[T,onlytextwidth]
    \begin{column}{.31\textwidth}
      \centering
      {\color{Muted}\fontsize{6.8}{7.5}\selectfont TEST ARCHITECTURE\par}\vspace{1mm}
      \paperimage{width=.94\linewidth}{assets/citation-1662-test-architecture.jpg}
    \end{column}
    \begin{column}{.31\textwidth}
      \centering
      {\color{Muted}\fontsize{6.8}{7.5}\selectfont FLATSAT\par}\vspace{1mm}
      \paperimage{width=.94\linewidth}{assets/citation-1662-flatsat.jpg}
    \end{column}
    \begin{column}{.31\textwidth}
      \centering
      {\color{Muted}\fontsize{6.8}{7.5}\selectfont COMPRESSED TEST SEQUENCE\par}\vspace{1mm}
      \paperimage{height=2.45cm}{assets/citation-1662-test-timeline.jpg}
    \end{column}
  \end{columns}
  \vspace{1mm}
  \begin{compactitems}
    \item Connection stages progress from direct computer commands, to a hardwired radio path, to the spacecraft and ground antenna path. The sequence compresses long operational intervals into an eight-hour campaign.
    \item Planned EPS checks include solar-panel illumination, resets and recovery, load shedding, and voltage, current, and temperature monitoring. The complete functional campaign remained pending.
  \end{compactitems}
  \vfill
  \source{Source figures: \emph{Abstracting CubeSat Operations: A Path to Real Cubesat Interoperability}, Figures 23--25, pp. 109--113.}
\end{frame}

% -----------------------------------------------------------------------------
\begin{frame}{Integrated EPS telemetry during the DESCENT campaign}
  \begin{columns}[T,onlytextwidth]
    \begin{column}{.58\textwidth}
      \centering
      \paperimage{width=.95\linewidth}{assets/citation-1662-telemetry-gui.jpg}
    \end{column}
    \begin{column}{.38\textwidth}
      \begin{tightitems}
        \item The telemetry view exposes battery current and voltage, plus panel current, voltage, and temperature.
        \item It records switch currents and states, EPS status and resets, and events, errors, and fatal state.
        \item Live and delayed data support operators during a test and preserve evidence for later analysis.
        \item The data connect a command or load change to the resulting EPS response.
      \end{tightitems}
    \end{column}
  \end{columns}
  \vfill
  \source{Source figure: \emph{Abstracting CubeSat Operations: A Path to Real Cubesat Interoperability}, Figure 26, pp. 114--120.}
\end{frame}

% -----------------------------------------------------------------------------
\begin{frame}{Model-based reliability analysis prioritizes EPS failure modes}
  \begin{columns}[T,onlytextwidth]
    \begin{column}{.53\textwidth}
      \centering
      \paperimage{width=.98\linewidth}{assets/review-1158-criticality-comparison.jpg}
    \end{column}
    \begin{column}{.43\textwidth}
      {\fontsize{8.2}{9.5}\selectfont
      \begin{compactitems}
        \item The SysML model integrates EPS architecture with FMEA, event scenarios, criticality analysis, and fault trees.
        \item Thirteen modes cover regulation, switching, protection, measurement, solar conversion, reverse current, bypass, and storage.
        \item Quantitative ranking placed battery storage and MPPT failure highest. The qualitative and quantitative rankings differed for some modes.
        \item The fault tree produced \metric{0.8802} reliability under simplifying assumptions, for comparison between designs.
      \end{compactitems}}
    \end{column}
  \end{columns}
  \vfill
  \source{Source: \emph{Model-Based Reliability Engineering: A methodology to enhance reliability analysis in student-based CubeSat design projects}, Figures 3.12 and 3.16, pp. 57--63.}
\end{frame}

% -----------------------------------------------------------------------------
\begin{frame}{Model-based diagnosis of a boost-converter fault before a hard limit}
  \begin{columns}[T,onlytextwidth]
    \begin{column}{.78\textwidth}
      \centering
      \paperimage{width=.98\linewidth}{assets/review-1251-boost-fault.jpg}
    \end{column}
    \begin{column}{.18\textwidth}
      \centering
      \paperimage{width=.96\linewidth}{assets/review-1251-boost-diagnosis.jpg}\par\vspace{1mm}
      {\color{Muted}\fontsize{6.6}{7.2}\selectfont ISOLATED DIAGNOSIS\par}
    \end{column}
  \end{columns}
  \vspace{1mm}
  \begin{compactitems}
    \item Qualification-model housekeeping data were modified with subtle faults and compared with simulated expected values.
    \item In the boost-converter case, expected current was \metric{0.18 A} and measured current was \metric{0.34 A}. The diagnosis isolated \term{Boostconverter2} before a hard telemetry limit was crossed.
    \item Five experiments covered single, cyclic, and indirectly observable faults. All injected discrepancies were detected and the reported root causes were identified.
  \end{compactitems}
  \vfill
  \source{Source figures and results: \emph{Model-Based Fault Detection and Diagnosis for Spacecraft with an Application for the SONATE Triple Cube Nano-Satellite}, Figures 20 and 22 and Table 2, pp. 20--31. Faults were injected into qualification-model data.}
\end{frame}

% -----------------------------------------------------------------------------
\begin{frame}{In-orbit schedule validation with GOMX-3 battery telemetry}
  \begin{columns}[T,onlytextwidth]
    \begin{column}{.48\textwidth}
      \centering
      \paperimage{width=.96\linewidth}{assets/review-1303-flight-schedule.jpg}\par\vspace{1mm}
      {\color{Muted}\fontsize{6.8}{7.5}\selectfont SCHEDULE, PREDICTION, AND TELEMETRY\par}
    \end{column}
    \begin{column}{.48\textwidth}
      \centering
      \paperimage{width=.96\linewidth}{assets/review-1303-voltage-load.jpg}\par\vspace{1mm}
      {\color{Muted}\fontsize{6.8}{7.5}\selectfont MEASURED LOADS AND STATE ESTIMATES\par}
    \end{column}
  \end{columns}
  \vspace{1mm}
  \begin{compactitems}
    \item The 24-hour run reached \metric{40.3\%} minimum state of charge, \metric{19.3\%} depletion risk, and two safe-mode entries. Later 36-hour and 60-hour runs had no safe-mode entry.
    \item Measured current and voltage exposed drift in the battery estimate, an initial net-power mismatch, and incorrect third-party power draws.
    \item Telemetry closed the loop from schedule execution to model correction and later model-predictive scheduling.
  \end{compactitems}
  \vfill
  \source{Source figures and data: \emph{Battery-aware scheduling in low orbit: the GomX--3 case}, Figures 9--10 and Table 1, pp. 13--22.}
\end{frame}

% -----------------------------------------------------------------------------
\begin{frame}{EPS sizing benchmark results: agreement and sensitivity}
  \begin{columns}[T,onlytextwidth]
    \begin{column}{.48\textwidth}
      \centering
      {\color{Muted}\fontsize{7}{7.8}\selectfont 1U REFERENCE\par}\vspace{1mm}
      {\fontsize{7.2}{8.5}\selectfont
      \begin{tabular}{l r}
        \hline
        Quantity & Signed error \\
        \hline
        Average power & 0.00\% \\
        Orbit energy & 0.66\% \\
        Battery capacity & 5.20\% \\
        Effective solar area & --7.40\% \\
        True solar area & --19.30\% \\
        \hline
      \end{tabular}}
    \end{column}
    \begin{column}{.48\textwidth}
      \centering
      {\color{Muted}\fontsize{7}{7.8}\selectfont POWERSAT REFERENCE\par}\vspace{1mm}
      {\fontsize{7.2}{8.5}\selectfont
      \begin{tabular}{l r}
        \hline
        Quantity & Signed error \\
        \hline
        Average power & 0.00\% \\
        Battery capacity & 19.02\% \\
        Battery-cell count & 50.00\% \\
        Average solar power & --1.73\% \\
        Effective solar area & 0.95\% \\
        \hline
      \end{tabular}}
    \end{column}
  \end{columns}
  \vspace{2mm}
  \begin{compactitems}
    \item Average power matched both references. The PowerSat solar metrics differed by less than 2\%.
    \item The 1U true solar-area result differed by \metric{--19.30\%}. PowerSat battery capacity differed by \metric{19.02\%}, and discrete cell count differed by \metric{50.00\%}.
    \item The benchmark supports selected calculations and exposes sensitivity to solar and battery assumptions.
  \end{compactitems}
  \vfill
  \source{Source: \emph{Development of a Cubesat Conceptual Design Tool and Implementation of the EPS Design Module}, Tables 16--17, pp. 89--91. Values are signed percent error against two reference designs.}
\end{frame}

\end{document}
